\emph{Setting and signs.} The background is the stored isotropic point: by Theorem \claimref{synthesis.scalar-kasner-wall-is-reduced-five-dimensional-vacuum-wall}(1) the five-dimensional vacuum points $(\tfrac12, \tfrac12, \tfrac12, -\tfrac12)$ and $(0, 0, 0, 1)$ both reduce to exponents $(\tfrac13, \tfrac13, \tfrac13)$ with $8\pi G q^{2} = \tfrac23$, and the pair solves the Einstein-scalar equations. The linearised computation is done in the component reading of the derived metric with the overall sign reversed, $-d\tau^{2} + \tau^{2/3}\sum_{i}(dx^{i})^{2}$, in which the equations read $\Ric = 8\pi G\,d\phi\otimes d\phi$; an overall sign of the metric changes no Christoffel symbol and no component of the Ricci tensor, so no mode equation depends on it. In that reading the perturbed components are $g_{00} = -(1 + 2f\cos kx^{1})$, $g_{01} = -b\sin kx^{1}$, $g_{ij} = \tau^{2/3}\delta_{ij}$; they are the first-order parts of the derived metric $\hb^{-1}(e_{\mu})\cdot\hb^{-1}(e_{\nu})$ of the stated field, since $\bigl((1 + F)e_{0}\bigr)^{2} = 1 + 2F + F^{2}$, $(1 + F)e_{0}\cdot(ae_{1} + Be_{0}) = B + FB$ and $(ae_{1} + Be_{0})^{2} = -a^{2} + B^{2}$. Write $E_{\mu\nu} = \delta\Ric_{\mu\nu} - 8\pi G\,\delta(\partial_{\mu}\phi\,\partial_{\nu}\phi)$, with $8\pi G\,\delta(\partial_{0}\phi\,\partial_{0}\phi) = \tfrac43\,\dot u\,\tau^{-1}\cos kx^{1}$ and $8\pi G\,\delta(\partial_{0}\phi\,\partial_{1}\phi) = -\tfrac23\,k\,u\,\tau^{-1}\sin kx^{1}$, all other components zero.

\emph{(i).} $E$ is linear in $(f, b, u)$, so it is the sum of three operators, each computed exactly from the Christoffel symbols and the Ricci tensor to first order. With $\beta$ the coefficient of $\sin kx^{1}$ in $g_{01}$ of the sign-reversed reading ($\beta = -b$):
\begin{align*}
E_{00} &= \bigl(-k^{2}\tau^{-2/3}f + \dot f/\tau + k\tau^{-2/3}\dot\beta - \tfrac43\dot u/\tau\bigr)\cos kx^{1},\\
E_{01} &= \tfrac23\,(k/\tau)\,(u - f)\,\sin kx^{1},\\
E_{11} &= \bigl(-\tfrac13\tau^{-1/3}\dot f + k^{2}f - k\dot\beta - \tfrac23(k/\tau)\beta\bigr)\cos kx^{1},\\
E_{22} = E_{33} &= \bigl(-\tfrac13\tau^{-1/3}\dot f - \tfrac13(k/\tau)\beta\bigr)\cos kx^{1},
\end{align*}
and $E_{02} = E_{03} = E_{12} = E_{13} = E_{23} = 0$ identically. For $k \neq 0$: $E_{01} = 0$ is $f = u$; $E_{22} = 0$ is $\beta = -\tau^{2/3}\dot f/k$, that is $b = \tau^{2/3}\dot u/k$. Substituting, $k\tau^{-2/3}\dot\beta = -\ddot u - \tfrac23\dot u/\tau$, so
\[
E_{00} = -\bigl(\ddot u + \dot u/\tau + k^{2}\tau^{-2/3}u\bigr)\cos kx^{1}, \qquad E_{11} = \tau^{2/3}\bigl(\ddot u + \dot u/\tau + k^{2}\tau^{-2/3}u\bigr)\cos kx^{1} .
\]
Hence the ten equations hold exactly when the three stated relations hold. The scalar's own equation is the contracted Bianchi identity of these and adds nothing. The background drops out of the equation for $u$: the first-derivative terms of $E_{00}$ sum to $\dot u/\tau - \tfrac23\dot u/\tau - \tfrac43\dot u/\tau = -\dot u/\tau$, the friction of a spectator, with no term in $u$ alone beside the gradient term; this is the statement that $a^{3}\dot\phi_{0}^{2}/H = 3q^{2}$ is constant at this point ($a = \tau^{1/3}$, $H = 1/(3\tau)$).

\emph{(ii).} With $s = \tfrac32k\tau^{2/3}$, $\dot u = k\tau^{-1/3}u_{s}$ and $\ddot u = k^{2}\tau^{-2/3}u_{ss} - \tfrac13k\tau^{-4/3}u_{s}$, so the equation is $k^{2}\tau^{-2/3}(u_{ss} + u_{s}/s + u) = 0$, using $k\tau^{-4/3}(1 - \tfrac13) = k^{2}\tau^{-2/3}/s$: Bessel's equation of order $0$, the stored equation at $p_{i} = \tfrac13$.

\emph{(iii).} Under $\tau \mapsto \tau + \xi^{0}$ at first order the scale factor in the spatial metric changes by $\tau^{2/3}(1 + 2H\xi^{0})$ and the background scalar by $\dot\phi_{0}\xi^{0}$, so with one sign convention for the induced change of the perturbations $\psi \mapsto \psi + H\xi^{0}$ and $\delta\phi \mapsto \delta\phi - \dot\phi_{0}\xi^{0}$, and $\psi + (H/\dot\phi_{0})\delta\phi$ is unchanged; a spatial displacement $\xi^{i} = \partial_{i}\chi$ changes the spatial metric by second derivatives of $\chi$ only and leaves $\delta\phi$ unchanged at first order, the background scalar being homogeneous. In the gauge of the statement $\psi = 0$, $\delta\phi = qu\cos kx^{1}$, $H/\dot\phi_{0} = (1/(3\tau))/(q/\tau) = 1/(3q)$, so $\zeta = \tfrac13u\cos kx^{1}$, and $f = u$ by (i). The pump field of $\zeta$ is $z = a\dot\phi_{0}/H = 3q\,a$.

\emph{(iv).} By (ii) the mode $u_{k}$ obeys the equation of Theorem \claimref{synthesis.kasner-approach-gives-horizon-exit-and-index-three}, whose clause (iii) gives, for the adiabatic vacuum with normalisation independent of $k$, $u_{k} \propto H^{(1)}_{0}(s)$ and the stated spectrum and index; $\zeta_{k} = u_{k}/3$ has the same. The cross modes at equal exponents have order $0$ by Theorem \claimref{synthesis.kasner-approach-coupled-modes-reach-measured-tilt-only-by-a-fitted-curvature-coupling}(v). That the vacuum normalisation of $u$ is independent of $k$ is the hypothesis of the clause; it is what the quadratic action of the curvature variable gives (textbook, not rederived here).
